\documentclass[12pt,a4paper]{article}
        \makeatletter
        \def\input@path{{../}}
        \makeatother
        \usepackage{almeria-fichas}

        \FichaMeta{Sesion 18}{Construimos la guia}{Bloque 4 · Producto final}{Integrar mapas, textos, tablas y graficas en un unico producto claro, ordenado y visualmente coherente.}{Procedimientos}{Aplicar, Analizar, Crear}

        \begin{document}

        \FichaCabecera

        \begin{materialesbox}
        \begin{itemize}
    \item Ficha impresa
    \item Lapiz
    \item Regla
    \item Materiales definitivos del equipo
    \item Ordenador o cartulina segun el formato elegido
\end{itemize}
        \end{materialesbox}

        \begin{pasosbox}
        \begin{enumerate}
    \item Reparte tareas claras dentro del equipo.
    \item Coloca primero el mapa y las secciones principales.
    \item Inserta despues tablas, calculos y graficas en el lugar adecuado.
    \item Revisa al final que todo tenga titulo, unidad y orden.
\end{enumerate}
        \end{pasosbox}

        \begin{recordatoriobox}
        \begin{itemize}
    \item Cada pagina o panel debe tener una funcion clara.
    \item Los datos matematicos tienen que verse y entenderse rapido.
    \item Un buen diseño deja espacio en blanco y evita saturar.
\end{itemize}
        \end{recordatoriobox}

        \begin{apoyobox}
        \begin{itemize}
    \item Antes de pegar o maquetar, haz un boceto rapido.
    \item Si una pagina tiene demasiado texto, resume.
    \item Usa el mismo estilo de titulos y colores en toda la guia.
\end{itemize}
        \end{apoyobox}

        \BloomSection{aplicar}

\BloomExercise{aplicar}{1}{Montaje guiado}{Organiza en una pagina estos elementos: mapa, tabla de monumentos, una grafica y un texto breve. Indica donde iria cada uno.}{8}

\BloomExercise{aplicar}{2}{Seccion util}{Diseña la distribucion de una seccion titulada \emph{Monumentos en cifras} con espacio para dato, unidad y explicacion.}{8}

\BloomSection{analizar}

\BloomExercise{analizar}{1}{Partes y relaciones}{Analiza un ejemplo relacionado con \emph{Construimos la guia}. Separa sus partes, indica cuales son relevantes y explica como se relacionan entre si.}{8}

\BloomExercise{analizar}{2}{Detecta errores o diferencias}{Compara dos respuestas, dos representaciones o dos procedimientos relacionados con esta sesion. Analiza sus diferencias y explica que errores o mejoras detectas.}{8}

\BloomSection{crear}

\BloomExercise{crear}{1}{Maqueta una pagina}{Crea el borrador de una pagina completa de la guia con titulo, subtitulo, mapa o imagen y al menos dos datos matematicos.}{10}

\BloomExercise{crear}{2}{Manual visual}{Crea una lista de tres normas de diseño para que todas las paginas del equipo mantengan el mismo estilo.}{7}

        \begin{evidenciabox}
        \begin{itemize}
    \item Organizar el producto final con orden y coherencia.
    \item Integrar calculos, imagenes y explicaciones.
    \item Tomar decisiones de diseño justificadas.
\end{itemize}
        \end{evidenciabox}

        \end{document}
